\subsubsection{Funktionsweise und Datenprodukte von LiDAR-Messungen}

Die Objekterfassung erfolgt durch das Aussenden eines modulierten Laserstrahls, der von der Oberfläche des Objekts reflektiert wird. Das System misst die Laufzeit (Time of Flight), die das Signal benötigt, um zum Empfänger zurückzukehren. Dieser fokussiert das reflektierte Licht auf einen Photodetektor, der die Photonen in elektrische Signale umwandelt. Ein Strahllenkungssystem steuert die Laserstrahlen in verschiedenen Winkeln, um das Sichtfeld (Field of View) abzutasten. Neben dem Laserscanner sind eine inertiale Messeinheit zur Messung von Eigenbewegung, Neigung und Orientierung des Sensors sowie ein globales Positionsbestimmungssystem zur zeitgenauen Erfassung der geografischen Position von Punkten integriert.

Dadurch ist es möglich pro Punkt die exakten X-Y-Z-Koordinaten im Raum zu definieren, die in ihrer Gesamtheit eine Punktwolke bilden. Neben der Position liefert das System Informationen über die Intensität der reflektierten Energie. Dies erlaubt Rückschlüsse auf das Oberflächenmaterial. Auch auf die Struktur der Oberfläche kann geschlossen werden, nämlich durch die Anzahl der Rückstrahlungen (Number of Returns) pro Strahl. Wenn ein Laserimpuls durch das Objekt dringt, erzeugen teildurchlässige Strukturen wie Vegetation mehrere Echos. Feste Oberflächen senden dagegen nur einen Rückstrahl. Ein Wert, der nicht mit dem LiDAR-System gemessen wird, ist der Normal Difference Vegetation Index (NDVI), der aus multispektralen Luftbildern berechnet wird. Ein Abgleich mit diesem Index ist jedoch sehr nützlich, da dieser als Indikator für die Vitalität von Vegetationsbedeckung fungiert.

Für die Analyse von Vegetation lassen sich insbesondere der höchste Rückstrahlpunkt (maximale Z-Koordinate) sowie die mittlere Höhe des Kronendachs und die Kronendichte ableiten. Die dreidimensionale Erfassung ermöglicht zudem die Bestimmung der vertikalen Vegetationsstruktur sowie die Ableitung von Geländestrukturen unter dem Blätterdach möglich.

  